\chapter{Conclusion and Future Work}
\label{ch:conclusion}

\section{Summary and Answers to the Research Questions}
\label{sec:concl-summary}

A monocular tennis line-calling system has to solve three perception sub-tasks, and this thesis studied each under one deterministic, like-for-like protocol. The hardest is tracking a small, fast ball that often spans only a handful of pixels. Locating the court is more forgiving, and it is what lets image coordinates be tied to the playing surface. The third sub-task, detecting when and where the ball bounces, is what ultimately drives the line call. Rather than proposing a single new architecture, the work framed each sub-task as a controlled comparison between candidate models trained and evaluated on a shared pipeline, and then composed the strongest candidate from each sub-task into one integrated program that produces a per-bounce in/out line call. The four research questions stated in Section~\ref{sec:intro-objectives} are answered in turn below; numbers are comparable only within a sub-task, because the ball and bounce experiments use the TrackNet tennis dataset while the court experiments use a separate court-keypoint dataset.

On ball tracking (RQ1), the comparison covered TrackNet V1--V5 together with the single-frame detector YOLO11m. TrackNetV4 was selected from the six-model subset comparison (Section~\ref{sec:ball-selection}) and then confirmed on the full split (Section~\ref{sec:ball-final}, Table~\ref{tab:ball-final}), where it reached a game10 test acc@5px of 0.956. It led on strict-tolerance localisation and on trajectory smoothness while remaining competitive in speed. The architectural progression from V1 to V5 was not monotone: the lightweight motion-attention of V4 outperformed the heavier temporal self-attention of V5, and the heatmap formulation shared by the TrackNet family clearly outperformed single-frame bounding-box regression, with YOLO11m the weakest localiser of the set.

On court detection (RQ2), four backbones were compared for 14-keypoint localisation. MobileNetV3-Small paired with a full-resolution decoder gave the best accuracy and latency trade-off (Section~\ref{sec:court-results}, Table~\ref{tab:court-results}): a PCK@7px of 0.992, a mean keypoint error of 2.46 px, a court-plane reprojection error of 7.67 cm, and a homography success rate of 0.999, all with only 1.28 million parameters and the fastest inference of the four. The smallest model was also the most accurate, and strict-tolerance accuracy was governed by the output heatmap resolution rather than by backbone capacity. The result is reliable enough to drive the homography on which the integrated system depends.

On bounce detection (RQ3), a gradient-boosted per-frame classifier (XGBoost) was compared with a temporal convolutional network (TCN) for recovering bounce timing from a noisy trajectory (Section~\ref{sec:bounce-results}, Table~\ref{tab:bounce-results}). Both solved the task well, each reaching an event F1 above 0.92 and almost never confusing a racket hit with a bounce, which a simple kinematic-reversal rule cannot do. They traded errors in a direction-dependent way: the TCN had the higher average precision (0.972 against 0.930), the better exact-frame timing (F1@k=0 of 0.717 against 0.680), and the higher recall, while XGBoost had the better tolerant event F1 (F1@k=3 of 0.948 against 0.925) and the higher precision (a single false positive across the match). The gap is small and rests on only 50 test events, so it should not be read as a decisive ranking. XGBoost was carried into the integrated system for its operating-point precision and its markedly simpler deployment.

On integration (RQ4), the three selected models composed into one end-to-end pipeline (Chapter~\ref{ch:integrated}). TrackNetV4, MobileNetV3, and XGBoost were wired into a single program that ingests a raw clip and emits an annotated video with a court minimap and, for each visible bounce, an in/out call obtained by projecting the bounce point through the homography and testing court-polygon containment (Section~\ref{sec:int-inout}, demonstrated in Section~\ref{sec:int-demo}).

\section{Future Work}
\label{sec:concl-future}

The most consequential gap concerns the line call itself. Because no dataset used in this work carries ground-truth in-or-out labels, the geometric verdict produced by the pipeline could not be scored against officiating-grade truth, and quantitative validation therefore requires collecting and annotating such labels before any accuracy claim can be made.

A second cluster of work concerns geometry and data. Monocular video cannot recover the true three-dimensional position of a bounce, so multi-camera fusion and explicit 3D trajectory reconstruction remain the principled path toward officiating accuracy. In parallel, larger and more varied datasets, spanning more matches, surfaces, camera angles, and broadcast styles, would address the domain-shift and small-test-set limitations of this study, most acutely the 50-event bounce test that underlies the RQ3 comparison. 

Finally, deployment and scope offer further work. The current CPU latency places the system in the offline regime, and real-time GPU serving would be needed to lift it toward live use. Beyond the three sub-tasks studied here, the perception stack could be extended to the player tracking, pose estimation, and shot classification that were deliberately placed out of scope (Section~\ref{sec:intro-scope}), broadening the system from line calling toward fuller match analysis.

\section{Closing Remarks}
\label{sec:concl-closing}

The benchmark isolates the design choices that matter within each sub-task, and the integrated demonstrator shows that the selected models compose into a coherent application. What is still missing is the quantitative-officiating layer: without ground-truth in/out labels the geometric verdict stays qualitative. The methodology and the working pipeline are meant to be reused and built on, not treated as a finished officiating system.
